% --- Archon physics formalization source begin ---
% archon:physics
% archon:covers PhyXMiniProblems/problem_phyx_mini_0159.lean
% archon:source-report reports/phyx_mini/problem_phyx_mini_0159.source.json

\chapter{Physics problem phyx\_mini\_0159}
\label{ch:PhyXMiniProblems_problem_phyx_mini_0159}

\paragraph{Problem source.}
\#\# Physical scenario

A \$4.0\textbackslash{};cm\$-diameter flower is \$200\textbackslash{};cm\$ from the \$50\textbackslash{};cm\$-focal-length lens of a camera.

\#\# Question

What is the diameter of the image on the detector?

\#\# Answer choices

- A: A: 1.5cm

- B: B: 1.3cm

- C: C: 1.7cm

- D: D: 1.9cm

\#\# Figure caption (auxiliary; use the image as primary evidence)

The image is an instructional diagram illustrating a ray tracing method for lenses. Here’s a detailed breakdown:

1. **Optical Axis**: A horizontal line across the image, marked as the optical axis, with labeled distances.

2. **Objects**:
   - A **green arrow** labeled \textbackslash{}( h \textbackslash{}) placed vertically on the left side of the lens, representing the object. It’s positioned on the optical axis.
   - A **lens** in the center, depicted as a blue symmetric convex shape.
   - On the right side, there’s an **inverted purple arrow** labeled \textbackslash{}( h' \textbackslash{}) representing the image.

3. **Focal Points**:
   - Two focal points, \textbackslash{}( f \textbackslash{}), are marked on either side of the lens along the optical axis. 

4. **Distances**:
   - The object distance is noted as \textbackslash{}( s = 200 \textbackslash{}, \textbackslash{}text\{cm\} \textbackslash{}).
   - The image distance is marked as \textbackslash{}( s' \textbackslash{}).
   - Distance from the object to the lens is labeled \textbackslash{}( 25 \textbackslash{}, \textbackslash{}text\{cm\} \textbackslash{}).

5. **Rays**:
   - Three purple lines are drawn indicating the light rays.
     - One ray is parallel to the optical axis before it refracts through the focal point on the right side.
     - Another ray passes through the near focal point on the left side and refracts parallel to the optical axis.
     - The third ray passes directly through the center of the lens without bending.

6. **Instructions**:
   - Step 1: "Lay out the optical axis, with a scale."
   - Step 2: "Draw the lens and mark its focal points."
   - Step 3: "Draw the object as an arrow with its base on the axis."
   - Step 4: "Draw the 3 special rays from the tip of the arrow." These rays are parallel to the axis, through the near focal point, and through the center of the lens.
   - Step 5: "The convergence point is the tip of the image. Draw the rest of the image."
   - Step 6: "Measure the image distance."

7. **Angles**:
   - Angles of refraction are labeled as \textbackslash{}( \textbackslash{}theta \textbackslash{}) at convergence points.

This diagram serves as a guide for visualizing how lenses form images using ray tracing.

\#\# Dataset metadata

Geometrical Optics; Spatial Relation Reasoning, Physical Model Grounding Reasoning

\paragraph{Recorded answer/context.}
B: B: 1.3cm

\paragraph{Figure/image path.}
/root/proposal\_for\_physic/science-mango/phyx\_mini\_run/phyx\_data/test\_image/159.png

\paragraph{Formalization target.}
create a compiling Lean file with sorry bodies at `PhyXMiniProblems/problem\_phyx\_mini\_0159.lean`. The Lean declarations must preserve the physical quantities, dimensions or dimensional roles, figure labels, governing-law hypotheses, and final relation expressed by this problem.
Use Mathlib/Physlib names found through LeanExplore where available. If a physics API is missing, introduce faithful local abstractions rather than scalar placeholder aliases.

\begin{theorem}[Physics formalization target]
\label{thm:physics:phyx_mini_0159:target}
\lean{PhyXMiniProblems.ProblemPhyXMini0159.problem_phyx_mini_0159}
\uses{def:physics:phyx-mini-0159:phyxminiproblems-problemphyxmini0159-lengthincentimeters, def:physics:phyx-mini-0159:phyxminiproblems-problemphyxmini0159-cameralenssetup, def:physics:phyx-mini-0159:phyxminiproblems-problemphyxmini0159-hasstatedcameradata, def:physics:phyx-mini-0159:phyxminiproblems-problemphyxmini0159-matchesprimaryfigure, def:physics:phyx-mini-0159:phyxminiproblems-problemphyxmini0159-usesrealimagedistanceconvention, def:physics:phyx-mini-0159:phyxminiproblems-problemphyxmini0159-hasphysicallengthconfiguration, def:physics:phyx-mini-0159:phyxminiproblems-problemphyxmini0159-obeysgaussianthinlensequation, def:physics:phyx-mini-0159:phyxminiproblems-problemphyxmini0159-obeystransversediameterlaw, def:physics:phyx-mini-0159:phyxminiproblems-problemphyxmini0159-iscorrectroundedanswer}
The assigned autoformalize agent should translate this physics problem into Lean declarations in the covered file, with theorem and lemma proof bodies written as `by sorry`.
\end{theorem}
\begin{proof}
This is an autoformalization task, not a proof task. Produce statements that can later be proved by the physics prover without weakening the physical model.
\end{proof}

% --- Archon declaration topology begin ---
\section{Declaration topology}
\begin{definition}[Optical Length]
  \label{def:physics:phyx-mini-0159:phyxminiproblems-problemphyxmini0159-opticallength}
  \lean{PhyXMiniProblems.ProblemPhyXMini0159.OpticalLength}
  A signed physical length whose scalar representation depends coherently on units
\end{definition}
\begin{proof}
  This is the declared model definition of Optical Length.
\end{proof}
\begin{definition}[centimeter Unit Choices]
  \label{def:physics:phyx-mini-0159:phyxminiproblems-problemphyxmini0159-centimeterunitchoices}
  \lean{PhyXMiniProblems.ProblemPhyXMini0159.centimeterUnitChoices}
  SI base units with centimeters selected as the length unit
\end{definition}
\begin{proof}
  This is the declared model definition of centimeter Unit Choices.
\end{proof}
\begin{definition}[length In Centimeters]
  \label{def:physics:phyx-mini-0159:phyxminiproblems-problemphyxmini0159-lengthincentimeters}
  \lean{PhyXMiniProblems.ProblemPhyXMini0159.lengthInCentimeters}
  \uses{def:physics:phyx-mini-0159:phyxminiproblems-problemphyxmini0159-opticallength, def:physics:phyx-mini-0159:phyxminiproblems-problemphyxmini0159-centimeterunitchoices}
  The signed scalar readout of an optical length in centimeters
\end{definition}
\begin{proof}
  This is the declared model definition of length In Centimeters.
\end{proof}
\begin{definition}[Thin Lens Kind]
  \label{def:physics:phyx-mini-0159:phyxminiproblems-problemphyxmini0159-thinlenskind}
  \lean{PhyXMiniProblems.ProblemPhyXMini0159.ThinLensKind}
  The optical power sign represented by the profile of a thin lens
\end{definition}
\begin{proof}
  This is the declared model definition of Thin Lens Kind.
\end{proof}
\begin{definition}[Image Nature]
  \label{def:physics:phyx-mini-0159:phyxminiproblems-problemphyxmini0159-imagenature}
  \lean{PhyXMiniProblems.ProblemPhyXMini0159.ImageNature}
  Whether the refracted rays themselves meet at the image
\end{definition}
\begin{proof}
  This is the declared model definition of Image Nature.
\end{proof}
\begin{definition}[Image Orientation]
  \label{def:physics:phyx-mini-0159:phyxminiproblems-problemphyxmini0159-imageorientation}
  \lean{PhyXMiniProblems.ProblemPhyXMini0159.ImageOrientation}
  Transverse orientation of the image relative to the object
\end{definition}
\begin{proof}
  This is the declared model definition of Image Orientation.
\end{proof}
\begin{definition}[Principal Axis Orientation]
  \label{def:physics:phyx-mini-0159:phyxminiproblems-problemphyxmini0159-principalaxisorientation}
  \lean{PhyXMiniProblems.ProblemPhyXMini0159.PrincipalAxisOrientation}
  Orientation of the principal optical axis in the source figure
\end{definition}
\begin{proof}
  This is the declared model definition of Principal Axis Orientation.
\end{proof}
\begin{definition}[Figure Point]
  \label{def:physics:phyx-mini-0159:phyxminiproblems-problemphyxmini0159-figurepoint}
  \lean{PhyXMiniProblems.ProblemPhyXMini0159.FigurePoint}
  Labeled axial points in the ray-tracing diagram
\end{definition}
\begin{proof}
  This is the declared model definition of Figure Point.
\end{proof}
\begin{definition}[Figure Ray]
  \label{def:physics:phyx-mini-0159:phyxminiproblems-problemphyxmini0159-figureray}
  \lean{PhyXMiniProblems.ProblemPhyXMini0159.FigureRay}
  Labels of the three special rays named in step 4 of the figure
\end{definition}
\begin{proof}
  This is the declared model definition of Figure Ray.
\end{proof}
\begin{definition}[Principal Ray Role]
  \label{def:physics:phyx-mini-0159:phyxminiproblems-problemphyxmini0159-principalrayrole}
  \lean{PhyXMiniProblems.ProblemPhyXMini0159.PrincipalRayRole}
  The standard geometrical-optics role assigned to a special ray
\end{definition}
\begin{proof}
  This is the declared model definition of Principal Ray Role.
\end{proof}
\begin{definition}[Figure Angle Mark]
  \label{def:physics:phyx-mini-0159:phyxminiproblems-problemphyxmini0159-figureanglemark}
  \lean{PhyXMiniProblems.ProblemPhyXMini0159.FigureAngleMark}
  The two occurrences of the angle label θ on the central ray
\end{definition}
\begin{proof}
  This is the declared model definition of Figure Angle Mark.
\end{proof}
\begin{definition}[Camera Lens Setup]
  \label{def:physics:phyx-mini-0159:phyxminiproblems-problemphyxmini0159-cameralenssetup}
  \lean{PhyXMiniProblems.ProblemPhyXMini0159.CameraLensSetup}
  \uses{def:physics:phyx-mini-0159:phyxminiproblems-problemphyxmini0159-opticallength, def:physics:phyx-mini-0159:phyxminiproblems-problemphyxmini0159-thinlenskind, def:physics:phyx-mini-0159:phyxminiproblems-problemphyxmini0159-imagenature, def:physics:phyx-mini-0159:phyxminiproblems-problemphyxmini0159-imageorientation, def:physics:phyx-mini-0159:phyxminiproblems-problemphyxmini0159-principalaxisorientation, def:physics:phyx-mini-0159:phyxminiproblems-problemphyxmini0159-figurepoint, def:physics:phyx-mini-0159:phyxminiproblems-problemphyxmini0159-figureray, def:physics:phyx-mini-0159:phyxminiproblems-problemphyxmini0159-principalrayrole, def:physics:phyx-mini-0159:phyxminiproblems-problemphyxmini0159-figureanglemark}
  The physical quantities and figure labels in the camera setup. The object and image diameters are positive transverse magnitudes. Axial positions are signed. objectDistance and imageDistance are positive distance magnitudes for the pictured real-image configuration. The marked angles are dimensionless real readouts in radians
\end{definition}
\begin{proof}
  This is the declared model definition of Camera Lens Setup.
\end{proof}
\begin{definition}[Has Stated Camera Data]
  \label{def:physics:phyx-mini-0159:phyxminiproblems-problemphyxmini0159-hasstatedcameradata}
  \lean{PhyXMiniProblems.ProblemPhyXMini0159.HasStatedCameraData}
  \uses{def:physics:phyx-mini-0159:phyxminiproblems-problemphyxmini0159-lengthincentimeters, def:physics:phyx-mini-0159:phyxminiproblems-problemphyxmini0159-cameralenssetup}
  Problem-statement data. These are the three input measurements: flower diameter 4.0 cm, object distance 200 cm, and focal length 50 cm. No image distance or image diameter is specified here
\end{definition}
\begin{proof}
  This is the declared model definition of Has Stated Camera Data.
\end{proof}
\begin{definition}[Matches Primary Figure]
  \label{def:physics:phyx-mini-0159:phyxminiproblems-problemphyxmini0159-matchesprimaryfigure}
  \lean{PhyXMiniProblems.ProblemPhyXMini0159.MatchesPrimaryFigure}
  \uses{def:physics:phyx-mini-0159:phyxminiproblems-problemphyxmini0159-lengthincentimeters, def:physics:phyx-mini-0159:phyxminiproblems-problemphyxmini0159-cameralenssetup}
  Readouts and qualitative geometry taken from the primary figure. The lens is at the origin; the object and focal points agree with the stated scale; and the real, inverted image lies beyond the far focal point. Its exact position and diameter remain deliberately unspecified
\end{definition}
\begin{proof}
  This is the declared model definition of Matches Primary Figure.
\end{proof}
\begin{definition}[Uses Real Image Distance Convention]
  \label{def:physics:phyx-mini-0159:phyxminiproblems-problemphyxmini0159-usesrealimagedistanceconvention}
  \lean{PhyXMiniProblems.ProblemPhyXMini0159.UsesRealImageDistanceConvention}
  \uses{def:physics:phyx-mini-0159:phyxminiproblems-problemphyxmini0159-cameralenssetup}
  Axial-distance bookkeeping for the sign convention shown in the figure. The incoming object and outgoing real image both have positive distance magnitudes, while the axial coordinates themselves are signed
\end{definition}
\begin{proof}
  This is the declared model definition of Uses Real Image Distance Convention.
\end{proof}
\begin{definition}[Has Physical Length Configuration]
  \label{def:physics:phyx-mini-0159:phyxminiproblems-problemphyxmini0159-hasphysicallengthconfiguration}
  \lean{PhyXMiniProblems.ProblemPhyXMini0159.HasPhysicalLengthConfiguration}
  \uses{def:physics:phyx-mini-0159:phyxminiproblems-problemphyxmini0159-lengthincentimeters, def:physics:phyx-mini-0159:phyxminiproblems-problemphyxmini0159-cameralenssetup}
  Positivity conditions for the physical magnitudes in this real-image setup
\end{definition}
\begin{proof}
  This is the declared model definition of Has Physical Length Configuration.
\end{proof}
\begin{definition}[Obeys Gaussian Thin Lens Equation]
  \label{def:physics:phyx-mini-0159:phyxminiproblems-problemphyxmini0159-obeysgaussianthinlensequation}
  \lean{PhyXMiniProblems.ProblemPhyXMini0159.ObeysGaussianThinLensEquation}
  \uses{def:physics:phyx-mini-0159:phyxminiproblems-problemphyxmini0159-cameralenssetup}
  The Gaussian thin-lens equation 1/f = 1/s + 1/s', written without division as f (s + s') = s s'. It is required in every unit system, so this is a law about dimensionful lengths rather than an equation about centimeter placeholders
\end{definition}
\begin{proof}
  This is the declared model definition of Obeys Gaussian Thin Lens Equation.
\end{proof}
\begin{definition}[Obeys Transverse Diameter Law]
  \label{def:physics:phyx-mini-0159:phyxminiproblems-problemphyxmini0159-obeystransversediameterlaw}
  \lean{PhyXMiniProblems.ProblemPhyXMini0159.ObeysTransverseDiameterLaw}
  \uses{def:physics:phyx-mini-0159:phyxminiproblems-problemphyxmini0159-cameralenssetup}
  The magnitude form of transverse magnification, image diameter / object diameter = image distance / object distance, written homogeneously without division. Image inversion is recorded separately by imageOrientation; diameters themselves are nonnegative magnitudes
\end{definition}
\begin{proof}
  This is the declared model definition of Obeys Transverse Diameter Law.
\end{proof}
\begin{definition}[Answer Choice]
  \label{def:physics:phyx-mini-0159:phyxminiproblems-problemphyxmini0159-answerchoice}
  \lean{PhyXMiniProblems.ProblemPhyXMini0159.AnswerChoice}
  Labels of the four multiple-choice answers
\end{definition}
\begin{proof}
  This is the declared model definition of Answer Choice.
\end{proof}
\begin{definition}[diameter In Centimeters]
  \label{def:physics:phyx-mini-0159:phyxminiproblems-problemphyxmini0159-answerchoice-diameterincentimeters}
  \lean{PhyXMiniProblems.ProblemPhyXMini0159.AnswerChoice.diameterInCentimeters}
  \uses{def:physics:phyx-mini-0159:phyxminiproblems-problemphyxmini0159-answerchoice}
  The image-diameter readout in centimeters printed beside each choice
\end{definition}
\begin{proof}
  This is the declared model definition of diameter In Centimeters.
\end{proof}
\begin{definition}[Is Nearest Tenth Readout]
  \label{def:physics:phyx-mini-0159:phyxminiproblems-problemphyxmini0159-isnearesttenthreadout}
  \lean{PhyXMiniProblems.ProblemPhyXMini0159.IsNearestTenthReadout}
  reported is a one-decimal-place readout within half a tenth of exact
\end{definition}
\begin{proof}
  This is the declared model definition of Is Nearest Tenth Readout.
\end{proof}
\begin{definition}[Is Correct Rounded Answer]
  \label{def:physics:phyx-mini-0159:phyxminiproblems-problemphyxmini0159-iscorrectroundedanswer}
  \lean{PhyXMiniProblems.ProblemPhyXMini0159.IsCorrectRoundedAnswer}
  \uses{def:physics:phyx-mini-0159:phyxminiproblems-problemphyxmini0159-lengthincentimeters, def:physics:phyx-mini-0159:phyxminiproblems-problemphyxmini0159-cameralenssetup, def:physics:phyx-mini-0159:phyxminiproblems-problemphyxmini0159-answerchoice, def:physics:phyx-mini-0159:phyxminiproblems-problemphyxmini0159-isnearesttenthreadout}
  A choice correctly reports the physical image diameter to the displayed precision
\end{definition}
\begin{proof}
  This is the declared model definition of Is Correct Rounded Answer.
\end{proof}
\begin{lemma}[image Distance In Centimeters eq two hundred div three]
  \label{lem:physics:phyx-mini-0159:phyxminiproblems-problemphyxmini0159-imagedistanceincentimeters-eq-two-hundred-div-three}
  \lean{PhyXMiniProblems.ProblemPhyXMini0159.imageDistanceInCentimeters_eq_two_hundred_div_three}
  \uses{def:physics:phyx-mini-0159:phyxminiproblems-problemphyxmini0159-lengthincentimeters, def:physics:phyx-mini-0159:phyxminiproblems-problemphyxmini0159-cameralenssetup, def:physics:phyx-mini-0159:phyxminiproblems-problemphyxmini0159-hasstatedcameradata, def:physics:phyx-mini-0159:phyxminiproblems-problemphyxmini0159-hasphysicallengthconfiguration, def:physics:phyx-mini-0159:phyxminiproblems-problemphyxmini0159-obeysgaussianthinlensequation}
  The thin-lens equation with f = 50 cm and s = 200 cm determines the real image distance s' = 200/3 cm
\end{lemma}
\begin{proof}
  The conclusion follows from the governing relations and figure readouts named in the statement of image Distance In Centimeters eq two hundred div three.
\end{proof}
\begin{lemma}[image Diameter In Centimeters eq four div three]
  \label{lem:physics:phyx-mini-0159:phyxminiproblems-problemphyxmini0159-imagediameterincentimeters-eq-four-div-three}
  \lean{PhyXMiniProblems.ProblemPhyXMini0159.imageDiameterInCentimeters_eq_four_div_three}
  \uses{def:physics:phyx-mini-0159:phyxminiproblems-problemphyxmini0159-lengthincentimeters, def:physics:phyx-mini-0159:phyxminiproblems-problemphyxmini0159-cameralenssetup, def:physics:phyx-mini-0159:phyxminiproblems-problemphyxmini0159-hasstatedcameradata, def:physics:phyx-mini-0159:phyxminiproblems-problemphyxmini0159-hasphysicallengthconfiguration, def:physics:phyx-mini-0159:phyxminiproblems-problemphyxmini0159-obeysgaussianthinlensequation, def:physics:phyx-mini-0159:phyxminiproblems-problemphyxmini0159-obeystransversediameterlaw}
  The transverse-diameter law then gives the exact detector-image diameter 4/3 cm from the 4 cm flower diameter
\end{lemma}
\begin{proof}
  The conclusion follows from the governing relations and figure readouts named in the statement of image Diameter In Centimeters eq four div three.
\end{proof}
% --- Archon declaration topology end ---
% --- Archon physics formalization source end ---
